% GENERATED FILE. Edit canonical lessons, then run npm run generate:books.
% canonical-source-hash: fnv1a64:5899b3d421b73d64
% canonical-lessons: RU-C234-podmetat, RU-C234-sazhat, RU-C234-polivat, RU-C234-rasti, RU-C234-kormit

\chapter{To sweep, To plant, To water, To grow, To feed}
\label{ch:ru-to-sweep-to-plant-to-water-to-grow-to-feed}
\hlchaptermodality{234}

\begin{chapteropening}
\textbf{By the end of this chapter:} \emph{I can say to sweep, to plant, to water, to grow and to feed.}

Complete the last lesson of chapter 234: I can say to sweep, to plant, to water, to grow and to feed.
\end{chapteropening}

\section[\texorpdfstring{\ru{подметать} (podmetát)}{подметать (podmetát)}]{\ru{подметать} (podmetát) — to sweep}

\label{lesson:RU-C234-podmetat}



\begin{quote}
Before the new one: say the Russian for to show, then the Russian for to bring.
\end{quote}

\subsection*{You'll want to know: \ru{подметать}}
\textbf{\ru{подметать}} (\emph{podmetát}) — \textquotedblleft{}to sweep\textquotedblright{}.

\begin{cousinweb}
\textbf{\ru{подметать}} is \emph{\ru{под}-}, \textquotedblleft{}under\textquotedblright{}, with \emph{\ru{мести}}, \emph{\ru{метать}}, \textquotedblleft{}to sweep\textquotedblright{}: to sweep up.
\end{cousinweb}

\begin{grammarlens}[title={What you've built}]
One more word for this chapter.
\end{grammarlens}

\subsection*{Your turn}
\begin{itemize}
\raggedright
  \item \emph{Say it:} \emph{podmetát}
  \item \emph{Say it:} \emph{podmetát}, once more
  \item \emph{Say it:} say \emph{podmetát}
  \item \emph{Recall:} say the Russian for to show, then the Russian for to bring, then say \emph{podmetát} again
\end{itemize}

\begin{tcolorbox}[breakable,colback=teal!4,colframe=teal!35!black,title={Before you move on}]
What does \ru{подметать} mean? (\textquotedblleft{}To sweep\textquotedblright{}.) Say it once more.
\end{tcolorbox}
\section[\texorpdfstring{\ru{сажать} (sazhát)}{сажать (sazhát)}]{\ru{сажать} (sazhát) — to plant}

\label{lesson:RU-C234-sazhat}



\begin{quote}
Before the new one: say the Russian for to bring, then the Russian for to sweep.
\end{quote}

\subsection*{You'll want to know: \ru{сажать}}
\textbf{\ru{сажать}} (\emph{sazhát}) — \textquotedblleft{}to plant\textquotedblright{}.

\begin{cousinweb}
\textbf{\ru{сажать}} goes with \emph{\ru{садить}}, \textquotedblleft{}to seat, to plant\textquotedblright{} (Vasmer), the \textquotedblleft{}make sit\textquotedblright{} form of \emph{\ru{сидеть}}, \textquotedblleft{}to sit\textquotedblright{}: a planted seedling is set to sit in the ground.
\end{cousinweb}

\begin{grammarlens}[title={What you've built}]
One more word for this chapter.
\end{grammarlens}

\subsection*{Your turn}
\begin{itemize}
\raggedright
  \item \emph{Say it:} \emph{sazhát}
  \item \emph{Say it:} \emph{sazhát}, once more
  \item \emph{Say it:} say \emph{sazhát}
  \item \emph{Recall:} say the Russian for to bring, then the Russian for to sweep, then say \emph{sazhát} again
\end{itemize}

\begin{tcolorbox}[breakable,colback=teal!4,colframe=teal!35!black,title={Before you move on}]
What does \ru{сажать} mean? (\textquotedblleft{}To plant\textquotedblright{}.) Say it once more.
\end{tcolorbox}
\section[\texorpdfstring{\ru{поливать} (polivát)}{поливать (polivát)}]{\ru{поливать} (polivát) — to water}

\label{lesson:RU-C234-polivat}



\begin{quote}
Before the new one: say the Russian for to sweep, then the Russian for to plant.
\end{quote}

\subsection*{You'll want to know: \ru{поливать}}
\textbf{\ru{поливать}} (\emph{polivát}) — \textquotedblleft{}to water\textquotedblright{}.

\begin{cousinweb}
\textbf{\ru{поливать}} is \emph{\ru{по}-} with \emph{-\ru{ливать}}, from \emph{\ru{лить}}, \textquotedblleft{}to pour\textquotedblright{}: to pour over. \emph{\ru{лить}} is a cousin of Lithuanian \emph{líeti}, \textquotedblleft{}to pour\textquotedblright{} (Vasmer).
\end{cousinweb}

\begin{grammarlens}[title={What you've built}]
One more word for this chapter.
\end{grammarlens}

\subsection*{Your turn}
\begin{itemize}
\raggedright
  \item \emph{Say it:} \emph{polivát}
  \item \emph{Say it:} \emph{polivát}, once more
  \item \emph{Say it:} say \emph{polivát}
  \item \emph{Recall:} say the Russian for to sweep, then the Russian for to plant, then say \emph{polivát} again
\end{itemize}

\begin{tcolorbox}[breakable,colback=teal!4,colframe=teal!35!black,title={Before you move on}]
What does \ru{поливать} mean? (\textquotedblleft{}To water\textquotedblright{}.) Say it once more.
\end{tcolorbox}
\section[\texorpdfstring{\ru{расти} (rastí)}{расти (rastí)}]{\ru{расти} (rastí) — to grow}

\label{lesson:RU-C234-rasti}



\begin{quote}
Before the new one: say the Russian for to plant, then the Russian for to water.
\end{quote}

\subsection*{You'll want to know: \ru{расти}}
\textbf{\ru{расти}} (\emph{rastí}) — \textquotedblleft{}to grow\textquotedblright{}.

\begin{cousinweb}
\textbf{\ru{расти}} $\leftarrow$ Proto-Slavic \emph{*orsti}, \textquotedblleft{}to grow\textquotedblright{} — inherited, but its \emph{\ru{а}} is Church Slavonic, says Vasmer, in place of native \emph{\ru{рости}}. Vasmer reports comparisons with Greek \emph{órnymi}, \textquotedblleft{}I stir\textquotedblright{}, and Latin \emph{orior}, \textquotedblleft{}I rise\textquotedblright{}.
\end{cousinweb}

\begin{grammarlens}[title={What you've built}]
One more word for this chapter.
\end{grammarlens}

\subsection*{Your turn}
\begin{itemize}
\raggedright
  \item \emph{Say it:} \emph{rastí}
  \item \emph{Say it:} \emph{rastí}, once more
  \item \emph{Say it:} say \emph{rastí}
  \item \emph{Recall:} say the Russian for to plant, then the Russian for to water, then say \emph{rastí} again
\end{itemize}

\begin{tcolorbox}[breakable,colback=teal!4,colframe=teal!35!black,title={Before you move on}]
What does \ru{расти} mean? (\textquotedblleft{}To grow\textquotedblright{}.) Say it once more.
\end{tcolorbox}
\section[\texorpdfstring{\ru{кормить} (kormít)}{кормить (kormít)}]{\ru{кормить} (kormít) — to feed}

\label{lesson:RU-C234-kormit}



\begin{quote}
Before the new one: say the Russian for to water, then the Russian for to grow.
\end{quote}

\subsection*{You'll want to know: \ru{кормить}}
\textbf{\ru{кормить}} (\emph{kormít}) — \textquotedblleft{}to feed\textquotedblright{}.

\begin{cousinweb}
\textbf{\ru{кормить}} comes from \emph{\ru{корм}}, \textquotedblleft{}fodder, food\textquotedblright{} (Vasmer), which Vasmer compares with Irish \emph{cuirm}, \textquotedblleft{}beer\textquotedblright{}, and Gaulish \emph{kourmi}, a kind of beer.
\end{cousinweb}

\begin{grammarlens}[title={What you've built}]
One more word for this chapter.
\end{grammarlens}

\subsection*{Your turn}
\begin{itemize}
\raggedright
  \item \emph{Say it:} \emph{kormít}
  \item \emph{Say it:} \emph{kormít}, once more
  \item \emph{Say it:} say \emph{kormít}
  \item \emph{Recall:} say the Russian for to water, then the Russian for to grow, then say \emph{kormít} again
\end{itemize}

\begin{tcolorbox}[breakable,colback=teal!4,colframe=teal!35!black,title={Before you move on}]
What does \ru{кормить} mean? (\textquotedblleft{}To feed\textquotedblright{}.) Say it once more.
\end{tcolorbox}
